\section{Pooling Algorithm}
\label{app:pool}

Our pooling algorithm is presented in Algorithm~\ref{alg:pool}.

\begin{figure}[h]
\centering
\fbox{\begin{minipage}{0.92\linewidth}
\textbf{Algorithm 1: Pooling}\\[2pt]
\textbf{Require:} Tensor $A \in \mathbb{R}^{L_M \times N \times D_M}$, target sizes $L_p, N_p$\\
\textbf{Ensure:} Tensor $A \in \mathbb{R}^{L_p \times N_p \times D_M}$\\
1: $A \leftarrow \mathrm{permute}(A, [D_M, L_M, N])$\\
2: Pad $A$ so that $L_M$ and $N$ are divisible by $L_p$ and $N_p$\\
3: Compute patch sizes: $f_L = \frac{L_{\text{pad}}}{L_p}$, $f_N = \frac{N_{\text{pad}}}{N_p}$\\
4: Apply pooling:\\
\hspace*{2em} $A \leftarrow \mathrm{F.max\_pool2d}(A, \text{kernel\_size} = (f_L, f_N))$\\
5: $A \leftarrow \mathrm{permute}(A, [L_p, N_p, D_M])$\\
6: \textbf{return} $A$
\end{minipage}}
\caption*{}
\label{alg:pool}
\end{figure}

\section{Permutation Symmetries}
\label{app:permutation}

Given the same number of layers and hidden dimensions, two LLMs can compute the exact same
function and yet produce entirely different activation tensors. This discrepancy stems from internal
symmetries in the model's weight space. To illustrate this, we consider a simplified transformer
architecture with $L$ layers, omitting residual connections and batch normalization for clarity (though
the analysis extends to those cases as well).

Recall that $A_l$ is the activation tensor at layer $l \in [L]$. Let $(Q_l, K_l, V_l)$ denote the query, key, and
value matrices at layer $l$, computed as:
\begin{equation}
(Q_l, K_l, V_l) = (A_{l-1} W_l^{Q_l}, A_{l-1} W_l^{K_l}, A_{l-1} W_l^{V_l}).
\end{equation}

Assume the feed-forward network (FFN) at layer $l$ produces output:
\begin{equation}
A_l = \mathrm{ReLU}(\mathrm{Attn}(A_{l-1}) W_l^1 + b_l^1) W_l^2 + b_l^2.
\label{eq:ffn}
\end{equation}

Now assume the model uses $d$-dimensional hidden features, and let $P \in \mathbb{R}^{d \times d}$ be a permutation
matrix, and define the following weight transformation:
\begin{align}
(W_l^2, b_l^2) &\to (W_l^2 P^\top, b_l^2 P^\top), \\
(W_{l+1}^{Q_{l+1}}, W_{l+1}^{K_{l+1}}, W_{l+1}^{V_{l+1}}) &\to (P W_{l+1}^{Q_{l+1}}, P W_{l+1}^{K_{l+1}}, P W_{l+1}^{V_{l+1}}).
\end{align}
Although the new weights (the right hand side) are very different, it is straightforward to verify that
the permutation matrices cancel out, leaving the underlying function unchanged. However, recalling
Equation~\eqref{eq:ffn}, the activation tensors differ significantly, as the transformation permutes the feature
dimension as follows:
\begin{equation}
A[l, n, d] \to A[l, n, \sigma(d)],
\end{equation}
where $\sigma$ is the permutation induced by $P$.

This illustrative example remarks the relevance of resorting to LLM-specific LA modules, as opposed to using a single shared linear adapter. Although padding can side-step the dimensionality
problem, such a single adapter would be required to implicitly learn (to be invariant to) all feature
permutations that may potentially arise ($D!$). While we note that principled approaches for handling
such symmetries exist -- such as designing layers that are invariant to feature permutations by design
-- we consider these less relevant to us given the limited number of considered LLMs. Exploring
symmetry-aware architectures still remains a promising direction we envision to explore in future
endeavors.

\paragraph{Extended Analysis to Standard Transformers Activations.} The implications extend naturally
to a standard transformer that includes residual connections and batch normalization. Residual
connections do not affect the analysis, since in our analysis the permutation symmetry is applies
uniformly across layers---the same permutation is used for all weight matrices, so the residual
addition remains consistent with our analysis. Batch normalization also preserves this symmetry, as
it is a permutation-equivariant operation: applying a permutation to the input features results in a
correspondingly permuted output.

\section{Broader Impact}

By enhancing the detection of hallucinations in Large Language Models, our work contributes to the
responsible development and deployment of generative AI by promoting greater transparency and
trust. However, we acknowledge that our findings also reveal aspects of the predictive information
embedded within LLM internals. This insight could be misused by malicious actors, potentially
motivating restrictions on LLM internals access and thereby slowing progress in open research.
